%-*-latex-*-

\documentclass[a4paper,11pt]{article}

% Language and fonts
%
\usepackage[utf8]{inputenc}
\usepackage[T1]{fontenc}
\usepackage[english]{babel}
\usepackage[hyphens]{url}
\usepackage{hyphenat}
\usepackage{array}
\usepackage{microtype}
\usepackage{ebgaramond}

\title{Notes from the Beginner Hungarian class at EKKE}
\date{July \oldstylenums{2026}}
\author{Christian Rinderknecht}

\begin{document}

\maketitle

\begin{itemize}

  \item \emph{név}: `name' (root).

  \item \emph{kereszt}: `cross' $\rightarrow$
    \emph{keresztneve\textbf{m}}: `my Christian name'.

  \item \emph{beceneve\textbf{m}}: `my nickname'.

  \item \emph{Mi a neve\textbf{d}?}: `What is your name?'; or
    \emph{Hogy hívnak?}. Answer: \emph{A nevem ...}; or
    \emph{Zoltán\textbf{nak} hív\textbf{nak}.}: `I am called Zoltán.'

  \item \emph{hívni}: `to call (someone by a name)'.

  \item \emph{melyik}: interr. `which (of several)'.

  \item \emph{vagy}: `or'; also `[you (sg.)] are' from \emph{lenni}:
    `to be'.

  \item \emph{szólíts}: `call me (imp.)'; e.g. \emph{Szólíts
  Zoli\textbf{nak}}.

  \item back vowels: \emph{autó} + \emph{á} + \emph{ú} + \emph{o}
    $\rightarrow$ endings \emph{-ak/-ok}, \emph{-nak}, \emph{-ok},
    \emph{-ul}.

  \item front vowels: \emph{teniszütő} (\emph{teniz} + \emph{ütő}
    `racket') + \emph{é} + \emph{í} + \emph{ű} + \emph{ö}
    $\rightarrow$ endings \emph{-nek}, \emph{-ek}, \emph{-ül}.

  \item mixed vowels: use back vowel endings, e.g. \emph{Ági\textbf{nak}}.


  \item \emph{légy szíves}: `please' (informal). \emph{szív}: `heart'.

  \item \emph{legy\textbf{en} szíves}: `please' (formal).

  \item \emph{kérem}: `I would like' (formal).

  \item  \emph{még egyszer}: `once more', `one more time'.

  \item \emph{írni}: `to write'.

  \item \emph{sajnos}: adj. `unfortunately'.

  \item \emph{értem}: `I understand'.

  \item \emph{kell}: `must, ought'.

  \item \emph{Ez jó?}: `[Is] this good/correct?'

  \item \emph{milyen}: interr. `what (kind of)'.

  \item Intonation in a question without marker (like \emph{milyen}):
    up on the penultimate syllable, down on the last.

  \item \emph{ki}: interr. `who'.

  \item \emph{Hogy vagy?}: `How are [you (sg.)]?'

  \item \emph{Hogy vagy\textbf{tok}?}: `How are [you (pl.)]?'

  \item \emph{Minden rendben?}: `[Is] everything all right?'

  \item \emph{jól}: adv. `well' (\emph{jó} + \emph{-l});
    \emph{Köszönöm, jól.}: `Thanks, I am fine.'; \emph{Nem túl jól.}:
    `Not too well.'

  \item \emph{sajnálom}: `I am sorry.' (formal)

  \item \emph{Egy kicsit meleg\textbf{em} van.}: `I'm a bit hot.'
    (Heat is possessed, hence the suffix \textbf{-em}). Otherwise:
    \emph{fáradt} (`tired'), \emph{álmos} (`sleepy'), \emph{éhes}
    (`hungry'), \emph{beteg} (`ill'), \emph{szomjas} (`thirsty').

  \item \emph{Fáj a tork\textbf{om}.}: `I have a sore throat.'

  \item \emph{Fáj a láb\textbf{am}.}: `My leg hurts.'

  \item \emph{örülök}: `I'm glad.' (informal); \emph{örvendek}
    (formal).

  \item \emph{Ki a barátod?}: `Who [is] your friend?'.

  \item \emph{Ez egy német autó.}: `It is \textbf{a} German car.' Here
    \emph{egy} is the indefinite article. Note: \emph{egy} functions
    in other contexts as the cardinal `one' (of many).

  \item \emph{Ez\textbf{ek} német autók.}: `Those are German cars.'
    Note the singular of \emph{német}.

  \item \emph{de}: conj. `but'

  \item \emph{Milyen nyelveken beszélsz?}: `What languages do you
    [sg.] speak?'; \emph{nyelv} + \emph{ek} (pl.) + \emph{-en} `in'

  \item \emph{(Én) Folyékonyan beszél\textbf{ek} spanyol\textbf{ul}.}:
    `I am fluent in Spanish.'; \emph{Egy kicsit beszélek koreaiul, de
  sajnos nem beszélek kínaiul.}: `I speak a little of Korean, but I
    unfortunately do not speak Chinese.'.

  \item Countries: \emph{Magyar\textbf{ország}},
    \emph{Német\textbf{ország}}, \emph{Francia\textbf{ország}},
    \emph{Anglia}, \emph{A\-me\-ri\-ka}, \emph{Portugália},
    \emph{Ausztria}.

  \item \emph{Hol élsz?}: `Where [do you (sg.)] live?';
    \emph{Eger\textbf{ben} élek.}: `[I] live in Eger.' Exceptions:
    \emph{Magyarország\textbf{on}}: `In Hungary.', and small islands,
    e.g. \emph{Malt\textbf{án}} (from \emph{Malta}).

  \item \emph{egészségedre}: \emph{egész} `whole' + \emph{-ség}
    `state' $\rightarrow$ `health' + \emph{-ed} `your (sg.)' +
    \emph{-re} `to'.

  \item \emph{Ő Péter.}: `[He is] Péter.' (Verb omitted in 3rd pers.).

  \item \emph{tanulni}: `to learn'

  \item \emph{most}: `now'

  \item Plurals: \emph{-k}, \emph{-ok}, \emph{-ek}, \emph{-ök},
    \emph{-ak}.

  \item Endings \emph{-unk/-ünk} depend on back/front vowels. Same
    with \emph{-tek/-tok} and \emph{-nek/-nak}. \emph{lenni}: `to
    be'; \emph{beszélni}: `to speak'; \emph{élni}: `to live';
    \emph{tanulni}: `to learn':
    \begin{center}
      \begin{tabular}{l|llll}
        & \textbf{lenni} & \textbf{beszélni} & \textbf{élni} &
        \textbf{tanulni} \\
        \hline
        én & vagyok   & beszélek  & élek  & tanulok \\
        te & vagy     & beszélsz  & élsz  & tanulsz \\
        ő  & (van)    & beszél    & él    & tanul \\
        mi & vagyunk  & beszélünk & élünk & tanulunk \\
        ti & vagytok  & beszéltek & éltek & tanultok \\
        ők & (vannak) & beszélnek & élnek & tanulnak
      \end{tabular}
      \end{center}

  \item \emph{Mi a telefonszámod?} (\emph{telefon} + \emph{szám}:
    `number' + \emph{-od} poss. 2nd pers. sg.): `What is your [sg.]
    phone number?'

  \item \emph{Mi a cím\textbf{ed}?}: `What is your [sg.] address?'

  \item \emph{Örülök, hogy megismertelek/találkoztunk!}

  \item \emph{Ismerkedjünk meg!}: `Let's get acquainted!'
    (\emph{ismer}: vb. `to know').

  \item \emph{anyanyelv}: mother tongue (\emph{anya} + \emph{nyelv}).

  \item \emph{hosszú}: adj. `long'. Anton. \emph{rövid}: `short'.

  \item \emph{Mit kérsz?}: `What would you [sg.] like?'. \emph{kérni}:
    `to want'.

  \item \emph{erős}: adj. `strong'. Anton. \emph{gyenge} `weak'.

  \item \emph{kis teje\textbf{t}}: No need to add \emph{-it} to
    \emph{kis}, as the accusative suffix is already added to
    \emph{tej} (`milk').

  \item \emph{\underline{Kicsit beszélek} franci\textbf{ául}.}: `I
    speak a bit of French'. \emph{Kicsit beszélek} forms a group.

  \item \emph{ásványvíz}: mineral water
    (\emph{víz}). Acc. \emph{ásványv\textbf{i}z\textbf{et}}.

  \item \emph{kéz}: `hand'. Acc. \emph{k\textbf{e}z\textbf{et}}.

  \item Digits:\\[2mm]
    \begin{tabular}{@{}rlrlrlrl@{}}
      0 & \emph{nulla}     & 10 & \emph{tíz}        & 20 & \emph{húsz} &
      30 & \emph{\textbf{harminc}} \\
      1 & \emph{egy}       & 11 & \emph{t\textbf{i}z\textbf{en}egy}   & 21 &
      \emph{huszonegy} &
      40 & \emph{negy\textbf{ven}}\\
      2 & \emph{két/kettő} & 12 & \emph{tizenkettő} & 22 &
      \emph{huszonkettő} &
      50 & \emph{öt\textbf{ven}}\\
      3 & \emph{három} & & & & & 60 & \emph{hat\textbf{van}}\\
      4 & \emph{négy} & & & & & 100 & \emph{száz}\\
      5 & \emph{öt} & & & & & 200 & \emph{kétszáz}\\
      6 & \emph{hat} & & & & & 1000 & \emph{ezer}\\
      7 & \emph{hét} \\
      8 & \emph{nyolc} \\
      9 & \emph{kilenc}
    \end{tabular}

  \item \emph{Két barát}: `Two friends.' (Not \emph{kettő}, as there
    is a noun.)

  \item \emph{Esik az eső?}: `Is it raining?' (\emph{esik}: `to fall',
    \emph{eső}: `rain').

  \item \emph{Süt a nap.}: `The sun is shining'.

  \item \emph{tó}: `lake', \emph{folyó}: `river', \emph{patak}
    `creek'; \emph{folyik}: `to flow'.

  \item \emph{de/hanem}: conj. `but'

  \item \emph{sok}: `many'. Anton. \emph{kevés}: `few'.

  \item \emph{pázsit}: `lawn, grass'.

  \item \emph{cipőfűző}: `shoelace'

  \item \emph{felvágott}: `cold cuts'

  \item Accusative of \emph{sajt} (`cheese') is \emph{sajt\textbf{ot}}.

  \item \emph{főzelék}: `vegetable'.

  \item Accusative singular. of \emph{leves} (`soup') is
    \emph{leves\textbf{t}}.

  \item Accusative of \emph{gyümölcs} (`fruit') is
    \emph{gyümölcs\textbf{öt}}.

  \item Accusative of \emph{} \emph{tej} (`milk') is
    \emph{tej\textbf{et}}.

  \item \emph{kolbász}: `sausage'

  \item \emph{szénsav\textbf{as}}:
    adj. `sparkling'. Anton. \emph{szénsav\textbf{mentes}} (`still').

  \item \emph{csapolt}: `on tap'; \emph{üveges} `bottled'.

  \item \emph{üdítő}: `soft drink'. \emph{Mennyibe kerül az üdítő?}
    $\rightarrow$ \emph{1\ 200 forint.}

  \item `How much (is it)?': \emph{Mennyi az üdítő?} or \emph{Mennyibe
  kerül?} $\rightarrow$ \emph{1\ 200 forint.} Answer: \emph{Tessék}:
    `Here it is (when serving someone)'; \emph{Ennyi?}  or
    \emph{Még/Más valamit?}: `Something else?'

  \item \emph{Mit adhatok?}: `What would you like?' (in shops);
    \emph{adni}: `to give'; \emph{Mit parancsol?}: `What would you
    like to order?' (in restaurants); \emph{rendel}: vb. `to order;
    \emph{rendel\textbf{és}}: (an) order (Do not use in person!)

  \item \emph{drága}: adj. `costly, dear'. Anton. \emph{olcsó}.

  \item \emph{büfé}: `buffet'.

  \item \emph{épület}: `(a) building'

  \item \emph{első}: adj. `first'; \emph{második}: adj. `second';
    \emph{harmadik}: adj. `third'.

  \item \emph{sarok}: `corner (of a street)'; \emph{sark\textbf{on}}: `on
    the corner (note the elided \emph{o})'.

  \item \emph{mellett}: `next to'. \emph{Ki ül mellett\textbf{ed}?}
    $\rightarrow$ \emph{János ül mellett\textbf{em}.}

  \item \emph{csemege}: `grocery'.

  \item \emph{Elnézést}: `Excuse-me.'; \emph{Bocsánat}: `Excuse-me.'
    (formal); \emph{Bocsi}: `Excuse-me.' (informal)

  \item \emph{aztán}: `then, after (some event)'. Long form:
    \emph{azután}, from \emph{az} + \emph{után}.

  \item \emph{jól}: adv. `well', from \emph{jó} `good' + \emph{-l}
    adverb of manner suffix. \emph{Megvagyok, köszönöm.}: Neutral
    answer to `Hogy vagy?'

  \item \emph{tej}: `milk'; \emph{tej\textbf{es}}: `milky';
    \emph{Kérek tejes kávét.}: `I want a coffee with milk.'

  \item \emph{narancs}: `orange' (fruit); \emph{mogyoró}: `hazelnut'
    (fruit); \emph{ribizli}: `redcurrant' (French `groseille');
    \emph{szőlő}: `grape'; \emph{málna}: `raspberry' (French:
    `framboise'); \emph{(görög)dinnye}: `watermelon' (from Greece);
    \emph{sárgadinnye}: `melon'; \emph{eper}: `strawberry' (Watch the
    elision of \emph{e} in the acc. sg. `epret'!); \emph{cseresznye}:
    `cherry'; \emph{meggy}: `sour cherry' (French `merise');
    \emph{ananász}: `pineapple'; \emph{alma}: `apple'; \emph{dió}:
    `walnut'; \emph{őszibarack}: `peach'; \emph{sárgabarack}:
    `apricot' (Watch acc. sg. \emph{sárgabarack\textbf{ot}!});
    \emph{körte}: `pear'; \emph{szilva}: `plum'; \emph{birs}: `quince'
    (French: `coing').

  \item \emph{narancs\textbf{lé}}: `orange juice'; \emph{lé}
    $\rightarrow$ acc. sg. \emph{levet}, pl. \emph{levek}.

  \item \emph{olíva} (pronounced like \emph{ol\textbf{i}va}): `olive';
    \emph{extraszűz olíva olaj}: `extra-virgin olive oil'.

  \item \emph{krumpli}: `potato' (higher register, restaurant:
    \emph{burgonya}); \emph{hagyma}: `onion';
    \emph{\textbf{fok}\-ha\-gy\-ma}: `garlic'; \emph{sárgarépa}:
    `carrot'; \emph{fehér répa}: `turnip'; \emph{karfiol}:
    `cauliflower', \emph{káposzta}: `cabbage', \emph{paradicsom}:
    `tomato'; \emph{saláta}: `salad'.

  \item \emph{vásárlás}: `shopping'.

  \item \emph{akkor}: then (`in this/that case').

  \item \emph{betűz}: vb. to spell; \emph{Betűzné, kérem?}:
    `Would you please spell it (for me)?' (formal)

  \item \emph{Kérsz még egy\textbf{et}?}: `Would you like one more?'

  \item \emph{száraz}: `dry' (e.g. wine). Anton. \emph{édes}
    (`sweet').

  \item \emph{étel}: `food'; \emph{ital}: `(a) drink'.

  \item \emph{fűszer}: `spice'; \emph{fűszer\textbf{es}}:
    adj. `spicy'; \emph{sós}: adj. `salty'; \emph{csípő\textbf{s}}:
    adj. `with chili (hot paprika)'; \emph{keserű}: adj. `bitter';
    \emph{savanyú}: adj. `sour'.

  \item \emph{néha}: adv. `sometimes'. Anton. \emph{soha}: `never'.

  \item \emph{szerint\textbf{em}}: `I think', `according to me';
    \emph{szerint\textbf{ed}}: `you think', `according to you'.

  \item \emph{például}: `for example'.

  \item \emph{ön}: `you' (formal sg.) $\rightarrow$ conjugate 3rd
    pers. sg; \emph{önök}: `you (formal pl.)' $\rightarrow$ conjugate
    3rd pers. pl.

  \item \emph{Igen, és \underline{németül is} beszélek.}

  \item \emph{Mi van \underline{a kert mellett}?} (postposition).

  \item \emph{Sok fa és virág.}: `(There are) Many trees and flowers.'
    (No nouns in plural, as the plural is implied by the adverb
    \emph{sok}.)

  \item \emph{pénz}: `money'.

  \item \emph{fagylalt}: `ice-cream'. (Short: \emph{fagyi}.)

  \item ---~\emph{Bocsánat, hol van piac?} ---~\emph{Itt egyenesen,
  aztán (az első/a második/a harmadik) utc\textbf{án}
  (jobbra/balra). Ott van a bolt mellett.} ---~\emph{Szívesen.}
    ---~\emph{Nincs mit}.

  \item \emph{előtt}: post. `in front of', `foremost', `before';
    \emph{egyenesen}: `straight (ahead)'; \emph{balra}: `left';
    \emph{jobbra}: `right'; \emph{erre}: `this way'; \emph{arra}:
    `that way'. \emph{Az autó előtt.}: `In front of the
    car'. \emph{Ebéd előtt.}: `Before lunch'.

  \item \emph{nehéz}: `difficult', `heavy'. Anton. \emph{könnyű}:
    `easy', `light'.

  \item \emph{ideges}: adj. `nervous'.

  \item \emph{érdekes}: adj. `interesting'.

  \item \emph{Magyarország} (country), \emph{magyar} (nationality),
    \emph{magyarul} (language). Same pattern: \emph{Németország},
    \emph{Spanyolország}, \emph{Olaszország}, \emph{Horvátország},
    \emph{Franciaország}, \emph{Gö\-rög\-or\-szág}.

  \item Countries and their nationalities: \emph{Szerbia}
    (\emph{szerb}), \emph{Bulgária} (\emph{bolgár}), \emph{Anglia}
    (\emph{angol}), \emph{Japán} (\emph{japán}), \emph{Kína}
    (\emph{kínai}), \emph{Amerika} (\emph{ameri\-kai}),
    \emph{Ausztria} (\emph{osztrák}).

  \item \emph{diák}: `student', `pupil'.


  \item Accusative singular: \emph{paprika} $\rightarrow$
    \emph{paprik\textbf{át}}; \emph{pizza} $\rightarrow$
    \emph{pizz\textbf{át}}; \emph{pisztácia} $\rightarrow$
    \emph{pisztáci\textbf{át}}; \emph{banán} $\rightarrow$
    \emph{banán\textbf{t}}; \emph{szalámi} $\rightarrow$
    \emph{szalámi\textbf{t}}; \emph{spaghetti} $\rightarrow$
    \emph{spaghetti\textbf{t}}; \emph{hotdog} $\rightarrow$
    \emph{hotdog\textbf{ot}}; \emph{hamburger} $\rightarrow$
    \emph{hamburger\textbf{t}}; \emph{joghurt} $\rightarrow$
    \emph{joghurt\textbf{ot}}; \emph{tej} $\rightarrow$ \emph{tejet};
    \emph{cukor} $\rightarrow$ \emph{cuk\textbf{rot}}; \emph{kávé}
    $\rightarrow$ \emph{kávé\textbf{t}}.

  \item Accusative + adjective: \emph{kolbász} $\rightarrow$
    \emph{kolbász\textbf{osat}} (\emph{-os} + \emph{-at});
    \emph{sonka} $\rightarrow$ \emph{sonká\textbf{sat}} (\emph{-s} +
    \emph{-at}); \emph{szalámi} $\rightarrow$
    \emph{szalámi\textbf{sat}} (\emph{-s} + \emph{-at}).

  \item \emph{kérni}: `to want'; ; \emph{szeretni}: `to love';
    \emph{jönni}: `to come'; \emph{dolgozik}: `to work':
    \begin{center}
      \begin{tabular}{l|llll}
        & \textbf{kérni} & \textbf{szeretni} & \textbf{jönni}
        & \textbf{dolgozik} \\
        \hline
        én & kérek  & szeretek  & jövök  & dolgoz\textbf{om}\\
        te & kérsz  & szeretsz  & jössz  & dolgoz\textbf{ol}\\
        ő  & kér    & szeret    & jön    & dolgoz\textbf{ik}\\
        mi & kérünk & szeretünk & jövünk & dolgoz\textbf{unk}\\
        ti & kértek & szerettek & jöttök & dolgoz\textbf{tok}\\
        ők & kérnek & szeretnek & jönnek & dolgoz\textbf{nak}
      \end{tabular}
    \end{center}

  \item \emph{inkább}: `instead'.

  \item \emph{cím}: `address'; \emph{Mi a címed?}: `What is your
    address?'

  \item \emph{piros}: `red' (inanimate objects); \emph{vörös}: `red'
    (living creatures, e.g. hair, wine, blood, emotional \& political
    contexts, dark red, e.g. \emph{vörös zászló}); \emph{zöld}:
    `green'; \emph{sötétzöld}: `dark green'; \emph{kék}: `blue';
    \emph{világ\textbf{os}kék}: `light blue' \emph{barna}: `brown';
    \emph{fehér}: `white'; \emph{sárga}: `yellow'; \emph{narancsárga}:
    `orange'; \emph{rózsaszín}: `pink'.

  \item \emph{van} at the start of a sentence means `There is'.

  \item \emph{hét}: `seven', `week'. \emph{Hét nap van egy héten.}:
    `There are seven days in a week.' \emph{A hét napjai.}
    (\emph{nap}: `day' (here, otherwise: `sun') + \emph{-ja} poss. +
    \emph{-i} pl.): `The days of the week.'

  \item \emph{hétfő}: `Monday', \emph{kedd}: `Tuesday', \emph{szerda}:
    `Wednesday', \emph{csütörtök}: `Thursday', \emph{péntek}:
    `Friday', \emph{szombat}: `Saturday', \emph{vasárnap}: `Sunday'.

  \item On a day of the week: \emph{hétfő\textbf{n}},
    \emph{kedd\textbf{en}}, \emph{szerd\textbf{án}},
    \emph{csütörtök\textbf{ön}}, \emph{péntek\textbf{en}},
    \emph{szombat\textbf{on}}, \emph{vasárnap} (no suffix!).

  \item \emph{csak}: adv. `only'.

  \item \emph{szerencsére}: adv. `happily', `fortunately'.

  \item \emph{Van egyetem a városodban?} (\emph{város} + 2nd
    pers. sg. poss. single\hyp{}poss. \emph{-od} + \emph{-ban}): `Is
    there a university in your town?'

  \item \emph{hogy}: Also means `that' (connector).

  \item \emph{szabadidő} (\emph{szabad} + \emph{idő}): `leisure time'.

  \item \emph{úszni}: `to swim' (but \emph{\textbf{u}szoda}: `swimming
    pool'); \emph{kávézni}: `to drink coffee'; \emph{futni}: `to run';
    \emph{olvasni}: `to read'; \emph{bulizni}: `to party';
    \emph{biciklizni}: `to ride a bicycle'; \emph{strandolni}: `to
    spend time at the beach'; \emph{táncolni}: `to dance';
    \emph{kirándolni}: `to hike', `to go on an excursion', `to go on a
    trip'; \emph{vacsorázni}: `to have dinner'; \emph{reggelni}: `to
    have breakfast'; \emph{ebédelni}: `to have lunch'; \emph{főzni}:
    `to cook'; \emph{zenét hallgatni}: `to listen to music';
    \emph{vásárolni}: `to shop'; \emph{netezni}: `to use the
    internet'; \emph{nézni}: `to watch', `to look'; \emph{csinálni}:
    `to do'; \emph{takarítani}: `to clean'.

  \item \emph{Szeretek zenét hallgatni.}: `I like to listen to music.'

  \item \emph{Te szeret\textbf{ed} az almát?}: `Do you like the
    apple?' Note the definite conjugation (\emph{-ed}), as there is a
    definite object.

  \item \emph{imádom} + object: `I love smth.'; \emph{utálom} +
    object: `I dislike smth.'

  \item \emph{Mit szeretsz csinálni szabadidő\textbf{dben}?}: `What do
    you like to do in your spare time?' (\emph{-d} +
    \emph{-ben}). \emph{Szabadidőmben imádok olvasni, de utálok
    futni.}: `In my leisure time, I love to read, but I hate jogging.'

  \item \emph{szín}: `colour'.

  \item \emph{kenyér}: `bread' $\rightarrow$
    acc. sg. \emph{keny\textbf{eret}}; \emph{zsemle/zsömle}: `bun';
    \emph{kifli}: `croissant'; \emph{lekvár}: `jam'; \emph{tojás}:
    `egg'; \emph{vaj}: `butter'; \emph{kalács}: `cake' $\rightarrow$
    acc. sg. \emph{kalács\textbf{ot}}.

  \item \emph{repülőtér}: `airport'; \emph{kórház}: `hospital';
    \emph{patika/gyógyszertár}: `pharmacy'; \emph{rendőrség}: `police
    (station)'; \emph{szálloda}: `hotel'; \emph{posta}: `post office';
    \emph{kollégium}: `dormitory (university)'; \emph{állomás}:
    `station'; \emph{megálló}: `stop'; \emph{buszmegálló}: `bus stop';
    \emph{buszállomás}: `bus station'; \emph{vasútállomás}: `(modest)
    train station'; \emph{pályaudvar}: `(grand) train station';
    \emph{végállomás}: `terminus'; \emph{iskola}: `school';
    \emph{benzinkút}: `petrol station' (\emph{benzin} + \emph{kút}:
    `well'); \emph{cukrászda}: `confectionery'; \emph{pékség}:
    `bakery'; \emph{kávézó}: `café'; \emph{piac}: `market';
    \emph{étterem} (\emph{ét}: `food' (see \emph{étel}) +
    \emph{terem}: `room, hall'); \emph{vár}: `castle'; \emph{bolt}:
    `shop'; \emph{templom}: `Catholic church'; \emph{egyetem}:
    `university'; \emph{könyvtár}: (\emph{könyv}: `book' + \emph{tár}:
    `storage'): `library'; \emph{könyv\textbf{es}bolt} (\emph{könyv} +
    \emph{-es} adj.-forming + \emph{bolt}): `bookstore';
    \emph{törökfürdő} (\emph{török}: `Turk(ish)' + \emph{fürdő}:
    `bath'): `Turkish bath'; \emph{uszoda}: `swimming pool' (see
    vb. \emph{úszni}); \emph{üzlet}: `business'.

  \item \emph{zászló}: `(country) flag'.

  \item \emph{nevezetes}: adj. `famous' + \emph{-ség}: `state (of
    being)' + \emph{-e} poss. + \emph{-i} pl. $\rightarrow$
    \emph{nevezetességei}. \emph{Milyen nevezetesség\textbf{ek}
    van\textbf{nak} a város\textbf{od}ban?}: `What landmarks are there
    in your town?' \emph{Hétfő a hét első napja.}: `Monday is the
    first day of the week.'

  \item \emph{holnap}: `tomorrow'; \emph{tegnap}: `yesterday';
    \emph{holnap\textbf{után}}: `after tomorrow'.

  \item \emph{oldal}: `side', `page'.

  \item \emph{óra}: `time, watch (to read to time), clock, class'.

  \item The time 3:15 can be expressed in two ways: \emph{Három óra
  tizenöt perc}, or \emph{negyed négy}: `A quarter into the fourth
    hour.'

  \item The time 1:30 is \emph{fél kettő}, with \emph{fél}: `half'.

  \item The time 3:45 is \emph{háromnegyed négy}: `Three quarters into
    the fourth hour.'

  \item The exact time requires \emph{óra}, e.g. 8:00 is \emph{nyolc
  \textbf{óra}}.

  \item The above times are ambiguous: we do not know if they refer to
    the afternoon. The times of the day are:
    \begin{center}
      \begin{tabular}{>{\itshape}ll}
        hajnal   & `dawn' \\
        reggel   & `morning' \\
        délelőtt & idem \\
        dél      & `noon' \\
        délután  & `afternoon' \\
        este     & `evening' \\
        éjszaka  & `night' \\
        ejfél    & `midnight'
      \end{tabular}
    \end{center}
    Use those qualifiers before the time of the day, e.g. \emph{Reggel
    tíz óra.}

  \item \emph{másodperc}: `second'.

  \item \emph{között}: `between', e.g. \emph{1 és 2 óra között.}:
    `Between 1 and 2.'

  \item \emph{Öt perc \textbf{múlva}}: `In five minutes.'; \emph{A
  bolt öt perc múlva zár.}: `The shop closes in five minutes.';
    \emph{Öt múlva találkozunk.}: `We meet in five minutes.';
    \emph{Tíz perc múlva ott vagyok/leszek.}: `I will be there in five
    minutes.'; \emph{A bolt öt perc múlva nyit.}: `The shop opens in
    five minutes.'

  \item When asking about a specific time, use the suffix \emph{-kor}
    on \emph{óra}, like so: \emph{Hány óra\textbf{kor} jön a vonat?}:
    `At what time does the train arrive?' Also in the answer, after
    the hour: \emph{Fél tizenkettő\textbf{kor}}, or only use
    \emph{óra}: \emph{Fél tizenkét \textbf{óra}.} Note \emph{két}
    instead of \emph{kettő}, as the former expects a noun, here
    \emph{óra}.

  \item Question words:
    \begin{center}
    \begin{tabular}{@{}>{\itshape}ll@{}}
      \textbf{Mi} a családneved? & `\textbf{What} is your surname?'\\
      \textbf{Mit} kérsz? & `\textbf{What} do you want?'\\
      \textbf{Milyen} nyelven beszélsz? & `\textbf{What}
      language do you speak?'\\
      & (Implicitly: Among others.) \\
      \textbf{Melyik} könyv? & `\textbf{Which} book?' \\
      \textbf{Hol/Merre} él ő? & `\textbf{Where} does he live?'\\
      \textbf{Hogy} kell írni a neved? & `\textbf{How} do you spell
      your name?' \\
      \textbf{Mennyibe kerül} a kávé? & `\textbf{How much} is a
      coffee?'\\
      \textbf{Hová/Hova} mész? & `\textbf{Where} are you going?'\\
      \textbf{Honnan} jöttél? & `What is your nationality?'\\
      & (lit. `\textbf{Where} did you come \textbf{from}?').\\
      \textbf{Ma/Mikor} épült? & `\textbf{When} was it built?'\\
      \textbf{Hány} óra van? & `\textbf{What} time is it?'\\
      \textbf{Mennyi} az idő? & idem. \\
      \textbf{Hánykor} nyit a piac? & `\textbf{At what time} does
      the market open?' \\
      \textbf{Hány órakor} nyit a piac? & idem \\
      \textbf{Miért} van itt? & `\textbf{Why} is it here?' \\
      \textbf{Ki} vagy? & `\textbf{Who} are you?' (informal sg.) \\
      \textbf{Kivel} mész moziba? & `\textbf{With whom} do you go to
      the cinema?'
    \end{tabular}
    \end{center}

  \item Possible answers to \emph{Hol?} (`where (to be)'):
    \begin{itemize}

      \item Inside non-official buildings or countries or some places
        with natural elements: suffixes \emph{-ban/-ben}, depending on
        the vowel harmony, e.g. \emph{a kert\textbf{ben}} (`At the
        garden'), \emph{az erdő\textbf{ben}} (`In the forest'),
        \emph{az uszodá\textbf{ban}} (`At the swimming pool'),
        \emph{Franciaország\textbf{ban}} (`In France'), \emph{a
        cukrászdá\textbf{ban}}, \emph{a postá\textbf{ban}} (`At the
        post office'). Note the vocalic changes from \textbf{a} to
        \textbf{á} before the suffix.

      \item On an open area or official buildings or Hungary or
        events: suffixes \emph{-on/-ön/-en}, depending on vowel
        harmony: \emph{Magyarország\textbf{on}} (`In Hungary'),
        \emph{a koncert\textbf{en}} (`At the concert'), \emph{a
        piac\textbf{on}} (`At the market'), \emph{az
        egyetem\textbf{en}} (`At the university'), \emph{a
        repülőtér\textbf{en}} (`At the airport'), \emph{a
        vas\-út\-állomás\textbf{on}} (`At the train station'),
        \emph{Mindig a számítógép\textbf{en} dolgozom.}: `I always
        work with (on) the computer.'

      \item At someone's place: suffixes \emph{-nál/-nél},
        e.g. \emph{Anná\textbf{nál}}: `At Anna's' (Note the vocalic
        change from \textbf{a} to \textbf{á} before the suffix.),
        \emph{fodrász\textbf{nál}}: `At the hairdresser'.

    \end{itemize}
    There are no rules for towns: they can either be \emph{-ban/-ben},
    like \emph{Eger\textbf{ben}} or \emph{-on/-en}, like
    \emph{Budapest\textbf{en}}.

  \item Possible answers to \emph{Hová/Hová?} (`where to')
    \begin{itemize}

      \item Suffixes \emph{-ba/-be} whenever \emph{-ban/-ben}
        (respectively) would apply to the question \emph{Hol?},
        e.g. \emph{a gyógyszertár\textbf{ba}} (`to the pharmacy').

      \item Suffixes \emph{-ra/re} whenever \emph{-on/-ön/-en} would
        apply to the question \emph{Hol?}, e.g. \emph{a
        tér\textbf{re}} (`to the square'), \emph{a
        repülőtér\textbf{re}} (`to the airport').

      \item Suffixes \emph{-hoz/-höz/-hez} whenever \emph{-nál/-nél}
        would apply to the question \emph{Hol?}

    \end{itemize}

  \item \emph{mindjárt}: adv. `soon' (\emph{dj} is pronounced as
    \emph{gy}).

  \item \emph{Jól aludtál?}: `Did you sleep well?'(\emph{dt} is
    pronounced as \emph{tt}).

  \item \emph{érdekes}: adj. `interesting'

  \item \emph{izgalmas}: adj. `exciting'

  \item \emph{szemüveg}: `(eye)glasses'; \emph{szemüveg\textbf{es}ek}:
    `(people) wearing glasses'.

  \item \emph{dobni}: `to throw, to beat a drum'.

  \item \emph{elkapni}: `to catch'

  \item \emph{A kertben sok fa van.}: `There are many trees in the
    garden.' No plural on \emph{fa} because of \emph{sok} (like number
    of things, e.g. \emph{két fa}). But: \emph{A kertben fá\textbf{k}
    van\textbf{nak}.}: `There are trees in the garden.'

  \item \emph{szökőkút} (\emph{szökő} + \emph{kút}): `fountain'.

  \item \emph{madár}: `bird' $\rightarrow$ pl. \emph{madarak}.

  \item \emph{vízesés}: `waterfall'.

  \item \emph{hegy}: `mountain'.

  \item \emph{vel\textbf{ed}}: adv. `with you (informal sg.)'.

  \item \emph{szó}: `word' $\rightarrow$ pl. \emph{szavak}.

  \item \emph{tó}: `lake' $\rightarrow$ pl. \emph{tavak},
    acc. \emph{tavat}.

  \item Adverbs of time: \emph{mindennap} (`every day'), \emph{mindig}
    (`always'), \emph{gyakran} (`often'), \emph{néha} (`sometimes'),
    \emph{ritkán} (`rarely'), \emph{sohasem} or \emph{soha nem}
    (`never').



  \item \emph{szokni}: `to usually do something'; \emph{Szoktál kávét
  inni reggel?}: `Do you usually drink coffee in the morning?';
    \emph{Én reggel szoktam futni.}: `I usually go running in the
    morning.' Note: The form \emph{szoktál} looks like the past of an
    irregular verb like \emph{enni} (see below), but this is a
    coincidence.


  \item \emph{kövér}: `fat'. Anton. \emph{sovány}: `slim'.

  \item \emph{ki-}: verbal prefix `out-'.

  \item \emph{naponta} (\emph{nap} + \emph{-on} + \emph{-ta}):
    adv. `daily'.

  \item `Neither ... nor...': \emph{Nem vagyok kövér, \underline{de
    nem} vagyok sovány \underline{sem}.} `I'm not fat, but I'm not
    thin either.'

  \item \emph{doboz}: `box'; also used as a classifier meaning `pack
    of'.

  \item Irregular verb \emph{lenni} `to be':
    \begin{center}
      \begin{tabular}{lll}
        Past & Present & Future\\
        \hline
        volt\textbf{am}   & vagyok           & lesz\textbf{ek} \\
        volt\textbf{ál}   & vagy             & lesz\textbf{el} \\
        volt              & van              & lesz \\
        volt\textbf{unk}  & vagy\textbf{unk} & lesz\textbf{ünk} \\
        volt\textbf{atok} & vagy\textbf{tok} & lesz\textbf{tek} \\
        volt\textbf{ak}   & van\textbf{nak}  & lesz\textbf{nek}
      \end{tabular}
    \end{center}

  \item \emph{már}: adv. `yet, already'; \emph{Voltál már Budapesten?}

  \item \emph{ma}: adv. `today' (or \emph{máma}: dialectal).

  \item \emph{biztos?}: `Really?'

  \item \emph{persze}: `of course'

  \item \emph{ének}, \emph{dal}: `song'; \emph{énekelni}: irreg. `to
    sing':
    \begin{center}
      \begin{tabular}{ll}
        & Énekelni\\
        \hline
        én & éne\textbf{klek} \\
        te & énekelsz \\
        ő  & énekel \\
        mi & ének\textbf{lünk} \\
        ti & énekeltek \\
        ők & énekelnek
      \end{tabular}
    \end{center}

  \item \emph{cuki}: adj. `cute'.

  \item \emph{híd}: `bridge', pl. \emph{hidak} (Not \emph{hidok}!),
    acc. sg. \emph{hidat}, acc. pl. \emph{hidakat}.

\end{itemize}


\end{document}
